\section{Especificação de Requisitos de Software}

Esta seção detalha as funcionalidades que o software deverá
oferecer e as condições que deverão ser atendidas durante sua
operação. Os requisitos são divididos em funcionais e não
funcionais.

\subsection{Requisitos Funcionais}

\subsubsection{RF01 -- Seleção de Ativos}

O software deverá permitir que o usuário selecione e remova os
ativos que participarão da análise da carteira.

\subsubsection{RF02 -- Entrada de Dados}

O software deverá permitir o carregamento de dados históricos pelo usuário e a obtenção de informações sobre os ativos por meio de APIs.

\subsubsection{RF03 -- Cálculo do Retorno Esperado}

O software deverá calcular o retorno esperado dos ativos e das
carteiras analisadas a partir dos dados fornecidos.

\subsubsection{RF04 -- Cálculo do Risco}

O software deverá calcular o risco das carteiras analisadas,
considerando a variância e a covariância entre os ativos que as
compõem.

\subsubsection{RF05 -- Geração de Carteiras}

O software deverá gerar diferentes combinações de pesos para os
ativos selecionados, respeitando as restrições estabelecidas pelo
modelo de Markowitz.

\subsubsection{RF06 -- Aplicação das Regras de Alocação}

O software deverá garantir que as carteiras analisadas respeitem
as regras de alocação definidas para o sistema, incluindo a soma
dos pesos igual a 1 e a ausência de pesos negativos quando não forem
permitidas vendas a descoberto.

\subsubsection{RF07 -- Análise de Risco e Retorno}

O software deverá calcular e apresentar a relação entre o risco e o
retorno esperado das carteiras geradas.

\subsubsection{RF08 -- Identificação da Fronteira Eficiente}

O software deverá identificar as carteiras que compõem a fronteira
eficiente segundo o critério estabelecido pelo modelo de Markowitz.

\subsubsection{RF09 -- Apresentação dos Resultados}

O software deverá apresentar os resultados da análise, incluindo os
pesos dos ativos, o risco e o retorno esperado das carteiras
selecionadas.

\subsubsection{RF10 -- Representação Gráfica}

O software deverá apresentar graficamente a relação entre risco e
retorno das carteiras analisadas, destacando a fronteira eficiente.

\subsection{Requisitos Não Funcionais}

\subsubsection{RNF01 -- Capacidade Operacional}

O software deverá permitir a análise de carteiras compostas por até
30 ativos simultaneamente.

\subsubsection{RNF02 -- Desempenho}

O software deverá realizar os cálculos necessários e apresentar a
representação gráfica dos resultados em, no máximo, 2 segundos para
uma análise contendo até 30 ativos.

\subsubsection{RNF03 -- Precisão Numérica}

O software deverá realizar os cálculos matemáticos com precisão
numérica adequada aos objetivos da análise de risco e retorno.

\subsubsection{RNF04 -- Validação dos Dados}

O software deverá validar os dados fornecidos antes da realização
dos cálculos e deverá informar o usuário quando houver dados
ausentes, inválidos ou insuficientes para a análise.

\subsubsection{RNF05 -- Consistência dos Resultados}

O software deverá utilizar os mesmos dados e critérios de cálculo
na apresentação dos resultados numéricos e na representação
gráfica.

\subsubsection{RNF06 -- Usabilidade}

O software deverá apresentar os dados de entrada e os resultados
da análise de forma clara e compreensível, permitindo que o usuário
utilize as principais funcionalidades sem necessidade de
conhecimento prévio sobre o funcionamento interno do sistema.
